\begin{abstract}
Many multiagent settings, such as poker, negotiation and security planning, require agents to
make decisions without knowing what their opponents know.
Strong strategies for these imperfect-information games can be learned through repeated
self-play, in which agents learn from their regret, but every round traverses the entire game
tree, so the number of iterations needed largely determines which games can be solved.
Counterfactual Regret Minimization (CFR) is the leading approach, and its variants speed it up
by changing how regrets are accumulated, with recent ones even adapting these choices
automatically.
However, these methods still apply the same update rule at every information set, even
though regret changes rapidly at some and has already stabilized at others.
In this work, we study two ways of adapting the update locally: \ours{}, in which each
information set chooses its own regret discount based on how volatile its recent regrets are,
and \mix{}, in which each information set learns online which of \cfrp{}, Discounted CFR and
Predictive \cfrp{} to follow, while remaining a no-regret method.
On five games, adapting only the discount gives modest gains, whereas \mix{}, without any
tuning, is the best or close to the best method on every game, while each existing variant
falls far behind on at least one.
\end{abstract}

\begin{IEEEkeywords}
Multiagent Systems, Game Theory, Counterfactual Regret Minimization, Imperfect-Information
Games, Nash Equilibrium
\end{IEEEkeywords}
